\chapter{2007年湖南卷（文）}

\section{单选题}

\begin{problem}
不等式 \(x^{2} > x\) 的解集是（\quad）

\choices
    {\((-\infty, 0)\)}
    {(0, 1)}
    {\((1, +\infty)\)}
    {\((-\infty, 0) \cup (1, +\infty)\)}
\end{problem}

\begin{problem}
若 \(O, E, F\) 是不共线的任意三点，则以下各式中成立的是（\quad）

\choices
    {\(\overrightarrow{EF} = \overrightarrow{OF} +\overrightarrow{OE}\)}
    {\(\overrightarrow{EF} = \overrightarrow{OF} -\overrightarrow{OE}\)}
    {\(\overrightarrow{EF} = -\overrightarrow{OF} +\overrightarrow{OE}\)}
    {\(\overrightarrow{EF} = -\overrightarrow{OF} -\overrightarrow{OE}\)}
\end{problem}

\begin{problem}
设 \( p: b^2 - 4ac > 0 (a \neq 0) \)，\( q \)：关于 \( x \) 的方程 \( ax^2 + bx + c = 0 (a \neq 0) \) 有实根，则 \( p \) 是 \( q \) 的（\quad）

\choices
    {充分不必要条件}
    {必要不充分条件}
    {充分必要条件}
    {既不充分也不必要条件}
\end{problem}

\begin{problem}
在等比数列 \(\{a_{n}\}\) （\(n \in \mathbf{N}^{*}\)）中，若 \(a_{1} = 1, a_{4} = \frac{1}{8}\)，则该数列的前 10 项和为（\quad）

\choices
    {\(2 - \frac{1}{2^8}\)}
    {\(2 - \frac{1}{2^9}\)}
    {\(2 - \frac{1}{2^{10}}\)}
    {\(2 - \frac{1}{2^{11}}\)}
\end{problem}

\begin{problem}
在 \((1 + x)^n\)（\(n \in \mathbf{N}^*\)）的二项展开式中，只有其 \(x^5\) 的系数最大，则 \(n =\)（\quad）

\choices
    {8}
    {9}
    {10}
    {11}
\end{problem}

\begin{problem}
如图，在正四棱柱 \(ABCD - A_{1}B_{1}C_{1}D_{1}\) 中，\(E\)、\(F\) 分别是 \(AB_{1}\)、\(BC_{1}\) 的中点，则以下结论中不成立的是（\quad）

\bitmapfigure[max width=0.38\linewidth,max totalheight=5.0cm]{img/2007/hunan_liberal/q6_fig1.png}

\choices
    {\(EF\) 与 \(BB_{1}\) 垂直}
    {\(EF\) 与 \(BD\) 垂直}
    {\(EF\) 与 \(CD\) 异面}
    {\(EF\) 与 \(A_{1}C_{1}\) 异面}
\end{problem}

\begin{problem}
根据某水文观测点的历史统计数据，得到某条河流水位的频率分布直方图（如图），从图中可以看出，该水文观测点平均至少一百年才遇到一次的洪水的最低水位是（\quad）

\bitmapfigure[max width=0.88\linewidth,max totalheight=6.8cm]{img/2007/hunan_liberal/q7_fig1.png}

\choices
    {48 米}
    {49 米}
    {50 米}
    {51 米}
\end{problem}

\begin{problem}
函数 \( f(x) = \begin{cases} 4x - 4, & x \leqslant 1, \\ x^2 - 4x + 3, & x > 1, \end{cases} \) 的图象和函数 \( g(x) = \log_2 x \) 的图象的交点个数是（\quad）

\choices
    {1}
    {2}
    {3}
    {4}
\end{problem}

\begin{problem}
设 \(F_{1}\)、\(F_{2}\) 分别是椭圆 \(\frac{x^{2}}{a^{2}} + \frac{y^{2}}{b^{2}} = 1 (a > b > 1)\) 的左、右焦点，\(P\) 是其右准线上纵坐标为 \(\sqrt{3} c\)（\(c\) 为半焦距）的点，且 \(|F_{1}F_{2}| = |F_{2}P|\)，则椭圆的离心率是（\quad）

\choices
    {\(\frac{\sqrt{3} - 1}{2}\)}
    {\(\frac{1}{2}\)}
    {\(\frac{\sqrt{5} - 1}{2}\)}
    {\(\frac{\sqrt{2}}{2}\)}
\end{problem}

\begin{problem}
设集合 \(M = \{1, 2, 3, 4, 5, 6\}\)，\(S_1, S_2, \cdots, S_k\) 都是 \(M\) 的含两个元素的子集，且满足：对任意的 \(S_i = \{a_i, b_i\}\)，\(S_j = \{a_j, b_j\}\)（\(i \neq j\)，\(i, j \in \{1, 2, 3, \cdots, k\}\)），都有 \(\min \left\{\frac{a_i}{b_i}, \frac{b_i}{a_i}\right\} \neq \min \left\{\frac{a_j}{b_j}, \frac{b_j}{a_j}\right\}\)（\(\min \{x, y\}\) 表示两个数 \(x, y\) 中的较小者）. 则 \(k\) 的最大值是（\quad）

\choices
    {10}
    {11}
    {12}
    {13}
\end{problem}

\section{填空题}

\begin{problem}
圆心为 \((1,1)\) 且与直线 \(x + y = 4\) 相切的圆的方程是\fillinblank{}.
\end{problem}

\begin{problem}
在 \(\triangle ABC\) 中，角 \(A \text{、} B \text{、} C\) 所对的边分别为 \(a \text{、} b \text{、} c\)，若 \(a = 1, c = \sqrt{3}\)，\(C = \frac{\pi}{3}\)，则 \(A =\)\fillinblank{}.
\end{problem}

\begin{problem}
若 \(a > 0, a^{\frac{3}{4}} = \frac{4}{9}\)，则 \(\log_{\frac{2}{3}} a =\)\fillinblank{}.
\end{problem}

\begin{problem}
设集合 \(A = \{(x, y) \mid y \geqslant |x - 2|, x \geqslant 0\}\)，\(B = \{(x, y) \mid y \leqslant -x + b\}\)，\(A \cap B \neq \emptyset\).

\begin{enumerate}[label={（\arabic*）},itemsep=0.2em]
\item \(b\) 的取值范围是\fillinblank{}；
\item 若 \((x, y) \in A \cap B\)，且 \(x + 2y\) 的最大值为 9，则 \(b\) 的值是\fillinblank{}.
\end{enumerate}
\end{problem}

\begin{problem}
棱长为 1 的正方体 \(ABCD - A_{1}B_{1}C_{1}D_{1}\) 的 8 个顶点都在球 \(O\) 的表面上，则球 \(O\) 的表面积是\fillinblank{}；设 \(E\)、\(F\) 分别是该正方体的棱 \(AA_{1}\)、\(DD_{1}\) 的中点，则直线 \(EF\) 被球 \(O\) 截得的线段长为\fillinblank{}.
\end{problem}

\section{解答题}

\begin{problem}
已知函数 \( f(x) = 1 - 2\sin^2\left(x + \frac{\pi}{8}\right) + 2\sin \left(x + \frac{\pi}{8}\right)\cos \left(x + \frac{\pi}{8}\right) \). 求：

\begin{enumerate}[label={（\arabic*）},itemsep=0.2em]
\item 函数 \( f(x) \) 的最小正周期；
\item 函数 \( f(x) \) 的单调增区间.
\end{enumerate}
\end{problem}

\begin{problem}
某地区为下岗人员免费提供财会和计算机培训，以提高下岗人员的再就业能力. 每名下岗人员可以选择参加一项培训、参加两项培训或不参加培训，已知参加过财会培训的有 60\%，参加过计算机培训的有 75\%，假设每个人对培训项目的选择是相互独立的，且各人的选择相互之间没有影响.

\begin{enumerate}[label={（\arabic*）},itemsep=0.2em]
\item 任选 1 名下岗人员，求该人参加过培训的概率；
\item 任选 3 名下岗人员，求这 3 人中至少有 2 人参加过培训的概率.
\end{enumerate}
\end{problem}

\begin{problem}
如图，已知直二面角 \(\alpha - PQ - \beta\)，\(A \in PQ\)，\(B \in \alpha\)，\(C \in \beta\)，\(CA = CB\)，\(\angle BAP = 45^{\circ}\)，直线 \(CA\) 和平面 \(\alpha\) 所成的角为 \(30^{\circ}\).

\begin{enumerate}[label={（\arabic*）},itemsep=0.2em]
\item 证明：\(BC \perp PQ\)；
\item 求二面角 \(B - AC - P\) 的大小.
\end{enumerate}

\bitmapfigure[max width=0.42\linewidth,max totalheight=4.8cm]{img/2007/hunan_liberal/q18_fig1.png}

\end{problem}

\begin{problem}
已知双曲线 \(x^{2} - y^{2} = 2\) 的右焦点为 \(F\)，过点 \(F\) 的动直线与双曲线相交于 \(A\)、\(B\) 两点，点 \(C\) 的坐标是 \((1, 0)\).

\begin{enumerate}[label={（\arabic*）},itemsep=0.2em]
\item 证明：\(\overrightarrow{CA} \cdot \overrightarrow{CB}\) 为常数；

\item 若动点 \(M\) 满足 \(\overrightarrow{CM} = \overrightarrow{CA} + \overrightarrow{CB} + \overrightarrow{CO}\)（其中 \(O\) 为坐标原点），求点 \(M\) 的轨迹方程.
\end{enumerate}
\end{problem}

\begin{problem}
设 \(S_{n}\) 是数列 \(\{a_{n}\} (n \in \mathbf{N}^{*})\) 的前 \(n\) 项和，\(a_{1} = a\)，且 \(S_{n}^{2} = 3n^{2}a_{n} + S_{n-1}^{2}\)，\(a_{n} \neq 0, n = 2, 3, 4, \cdots\).

\begin{enumerate}[label={（\arabic*）},itemsep=0.2em]
\item 证明：数列 \(\{a_{n+2} - a_n\} (n \geqslant 2)\) 是常数数列；
\item 试找出一个奇数 \(a\)，使以 18 为首项，7 为公比的等比数列 \(\{b_n\} (n \in \mathbf{N}^*)\) 中的所有项都是数列 \(\{a_n\}\) 中的项，并指出 \(b_n\) 是数列 \(\{a_n\}\) 中的第几项.
\end{enumerate}
\end{problem}

\begin{problem}
已知函数 \( f(x) = \frac{1}{3} x^3 + \frac{1}{2} ax^2 + bx \) 在区间 \([-1, 1), (1, 3]\) 内各有一个极值点.

\begin{enumerate}[label={（\arabic*）},itemsep=0.2em]
\item 求 \(a^2 - 4b\) 的最大值；
\item 当 \(a^2 - 4b = 8\) 时，设函数 \(y = f(x)\) 在点 \(A(1, f(1))\) 处的切线为 \(l\). 若 \(l\) 在点 \(A\) 处穿过 \(y = f(x)\) 的图象（即动点在点 \(A\) 附近沿曲线 \(y = f(x)\) 运动，经过点 \(A\) 时，从 \(l\) 的一侧进入另一侧），求函数 \(f(x)\) 的表达式.
\end{enumerate}
\end{problem}
